\subsection{Nuclear maps}

In this section, F denotes a Hilbert space over K. When K = C, we recall (I, p. 139, def. 3) that for \(u \in \mathcal{L}(\mathrm{E}; \mathrm{F})\) , we denote by \(|u|\) the positive endomorphism \(\sqrt{u^{*} \circ u}\) of E. When K = R, the element \(|u_{(\mathbf{C})}|\) of \(\mathcal{L}(\mathrm{E}_{(\mathbf{C})})\) is of the form \(v_{(\mathbf{C})}\) for a unique endomorphism \(v \in \mathcal{L}(\mathrm{E})\) , which is still denoted \(|u|\) (cf. I, p. 87).

We denote by \((u,v)\mapsto\langle u|v\rangle=\mathrm{Tr}(u^{*}v)\) the inner product in the Hilbert space \(\mathcal{L}_{2}(\mathrm{E};\mathrm{F})\) (EVT, V, p. 53, remarks 1 and 2).

For every \(u \in \mathcal{L}(\mathrm{E};\mathrm{F})\) , we denote by \(\| u\| _1\) the quantity \(\mathrm{Tr}(|u|)\) if \(|u|\) has finite trace and \(\| u\| _1 = +\infty\) if this is not the case. We therefore have \(\| u\| _1 \in \mathbf{R}_+ \cup \{+\infty\}\) . Since \(\| u\| = \| |u|\|\) (prop. 10, a) of I, p. 139) and since, if \(v\) is positive and of finite trace, then \(\mathrm{Tr}(v) \geqslant \| v\|\) (EVT, V, p. 49, formula (24bis) and p. 44, prop. 9), it follows that

\[
\| u \| \leqslant \| u \| _ {1}.\tag{9}
\]

If \(\| u\| _1\) is finite, then \(u\) is a Hilbert--Schmidt operator and \(\| u\| _2 \leqslant \| u\| _1\) (cor. 4 of IV, p. 165).

PROPOSITION 14. --- Let \(u \in \mathcal{L}_2(\mathrm{E};\mathrm{F})\) . We have

\[
\| u\|_{1} = \sup_{\substack{v\in \mathcal{L}_{2}(\mathrm{E};\mathrm{F})\\ \| v\| \leqslant 1}}|\langle v  |  u\rangle |,
\]

the supremum being computed in \(\mathbf{R}_{+} \cup \{+\infty\}\) .

Let \(\mathrm{B} = (e_i)_{i \in \mathrm{I}}\) be an orthonormal basis of the initial space \(\operatorname{Ker}(u)^{\circ}\) of u, \(\mathrm{C} = (f_i)_{i \in \mathrm{I}}\) be an orthonormal family of F and \((\alpha_i)_{i \in \mathrm{I}}\) be a family in \((\mathbf{R}_+^*)^{\mathrm{I}}\) such that \(u(e_i) = \alpha_i f_i\) for all \(i \in I\) (cor. 2 of IV, p. 150). The family \((\alpha_i)_{i \in \mathrm{I}}\) is square-integrable since u belongs to \(\mathcal{L}_2(\mathrm{E};\mathrm{F})\) (cor. 4 of IV, p. 165). Let \((e_j)_{j \in \mathrm{J}}\) be an orthonormal basis of E extending \((e_i)_{i \in \mathrm{I}}\) .

Let L be a finite subset of I. Let \(v_{L}\) be the continuous finite-rank linear map from E into F defined by

\[
v _ {\mathrm{L}} (x) = \sum_ {i \in \mathrm{L}} \langle e _ {i} | x \rangle f _ {i}
\]

for all \(x\) in E. We have \(\| v_{\mathrm{L}}\| \leqslant 1\) and \(v_{\mathrm{L}}\in \mathcal{L}_{2}(\mathrm{E};\mathrm{F})\) . Moreover, for every \(j\in \mathrm{J}\) , we have \(v_{\mathrm{L}}(e_j) = 0\) if \(j\notin \mathrm{L}\) and \(v_{\mathrm{L}}(e_j) = f_j\) if \(j\in \mathrm{L}\). Thus

\[
| \langle v _ {\mathrm{L}} | u \rangle | = | \operatorname{Tr} (v _ {\mathrm{L}} ^ {*} u) | = \left| \sum_ {j \in \mathrm{J}} \langle v _ {\mathrm{L}} (e _ {j}) | u (e _ {j}) \rangle \right| = \sum_ {j \in \mathrm{L}} \alpha_ {j}.
\]

From loc. cit., we deduce that

\[
\| u\|_{1} = \sup_{\text{L}}\sum_{j\in \text{L}}\alpha_{j}\leqslant \sup_{\substack{v\in \mathcal{L}_{2}(\text{E};\text{F})\\ \| v\| \leqslant 1}}|\langle v\mid u\rangle |\tag{10}
\]

in \(\mathbf{R}_{+}\cup\{+\infty\}\) .

This implies the equality of the proposition when \(\|u\|_{1}=+\infty\) .

Suppose that \(\|u\|_{1}\) is finite. For every \(i \in I\) , let \(p_{i}\) be the linear map from E to F such that \(p_{i}(e_{i}) = f_{i}\) and \(p_{i}(x) = 0\) for all \(x \in e_{i}^{\circ}\) . For all j and k in I, we have

\[
\langle p _ {j} \mid p _ {k} \rangle = \mathrm{Tr} (p _ {j} ^ {*} p _ {k}) = \sum_ {i \in \mathrm{J}} \langle p _ {j} (e _ {i}) \mid p _ {k} (e _ {i}) \rangle ,
\]

which is zero unless \(j = k\) , in which case this quantity equals \(\| f_k\|^2 = 1\) . The family \((p_i)_{i \in \mathrm{I}}\) is therefore orthonormal in \(\mathcal{L}_2(\mathrm{E};\mathrm{F})\) . Consequently, the family \((\alpha_i p_i)_{i \in \mathrm{I}}\) is summable in \(\mathcal{L}_2(\mathrm{E};\mathrm{F})\) ; its sum is equal to \(u\) since these two continuous linear maps coincide on the elements \(e_i\) for all \(i \in \mathrm{J}\) .

Let \(v \in \mathcal{L}_2(\mathrm{E};\mathrm{F})\) . For every \(i \in \mathrm{I}\) , it follows that

\[
\langle v \mid p _ {i} \rangle = \operatorname{Tr} (v ^ {*} p _ {i}) = \sum_ {j \in \mathrm{J}} \langle v (e _ {j}) \mid p _ {i} (e _ {j}) \rangle = \langle v (e _ {i}) \mid f _ {i} \rangle .
\]

If \(\|v\| \leqslant 1\) , we have

\[
\begin{array}{c} | \langle v | u \rangle | = \left| \langle v | \sum_ {i \in I} \alpha_ {i} p _ {i} \rangle \right| \leqslant \sum_ {i \in I} \alpha_ {i} | \langle v | p _ {i} \rangle | \\ = \sum_ {i \in I} \alpha_ {i} | \langle v (e _ {i}) | f _ {i} \rangle | \leqslant \sum_ {i \in I} \alpha_ {i} = \mathrm{Tr} (| u |), \end{array}
\]

whence the inequality

\[
\sup_{\substack{v\in \mathcal{L}_{2}(\mathrm{E};\mathrm{F})\\ \| v\| \leqslant 1}}|\langle v\mid u\rangle |\leqslant \| u\|_{1},
\]

which, combined with formula (10), concludes the proof of the proposition.

It follows from this proposition that the set \(\mathcal{L}_1(\mathrm{E};\mathrm{F})\) of continuous linear maps \(u\) from E to F such that \(\| u\| _1\) is finite is a vector subspace of \(\mathcal{L}_{2}(\mathrm{E};\mathrm{F})\) and that the map \(u\mapsto \| u\| _1\) is a semi-norm on \(\mathcal{L}_1(\mathrm{E};\mathrm{F})\) ; inequality (9) shows that it is a norm.

DEFINITION 4. --- We say that the vector space \(\mathcal{L}_{1}(\mathrm{E};\mathrm{F})\) equipped with the norm \(u\mapsto \| u\|_{1}\) is the space of nuclear maps from E to F. If \(u\in \mathcal{L}_1(\mathrm{E};\mathrm{F})\) , one says that \(u\) is nuclear.

Remarks. --- 1) Suppose that K = R. Let \(u \in \mathcal{L}(E; F)\) . We have \(u \in \mathcal{L}_{1}(E; F)\) if and only if \(u_{(\mathbf{C})} \in \mathcal{L}_{1}(E_{(\mathbf{C})}; F_{(\mathbf{C})})\) ; in this case, we have \(\|u\|_{1} = \|u_{(\mathbf{C})}\|_{1}\) .

\begin{enumerate}
\def\labelenumi{\arabic{enumi})}
\setcounter{enumi}{1}
\tightlist
\item
  Let \(u\) be a nuclear map from E to F. The map \(u\) being Hilbert--Schmidt, it is compact (cor. 2 of IV, p. 165). Moreover, proposition 14 implies that for every \(v \in \mathcal{L}_2(\mathrm{E};\mathrm{F})\) , we have
\end{enumerate}

\[
| \langle v | u \rangle | \leqslant \| v \| \| u \| _ {1}.\tag{11}
\]

\begin{enumerate}
\def\labelenumi{\arabic{enumi})}
\setcounter{enumi}{2}
\item
  Suppose that \(\mathrm{K} = \mathbf{C}\). Let \(u \in \mathcal{L}_1(\mathrm{E};\mathrm{E})\) such that \(u\) is positive. The norm \(\| u\| _1\) of \(u\) is the sum of the sequence \((\lambda_n(u))_{n\in \mathrm{I_E}}\) (cor. 3 of IV, p. 165).
\item
  The canonical inclusion of \(\mathcal{L}_1(\mathrm{E};\mathrm{F})\) in \(\mathcal{L}(\mathrm{E};\mathrm{F})\) has norm \(\leqslant 1\) (inequality (9), p. 167).
\end{enumerate}

PROPOSITION 15. --- Let G be a Hilbert space over K. The map \((u,v)\mapsto v\circ u\) from \(\mathcal{L}(\mathrm{E};\mathrm{F})\times\mathcal{L}(\mathrm{F};\mathrm{G})\) into \(\mathcal{L}(\mathrm{E};\mathrm{G})\) defines by passage to quotients a continuous bilinear map of norm \(\leqslant1\) from \(\mathcal{L}_{2}(\mathrm{E};\mathrm{F})\times\mathcal{L}_{2}(\mathrm{F};\mathrm{G})\) into \(\mathcal{L}_{1}(\mathrm{E};\mathrm{G})\) .

Let \(u \in \mathcal{L}_2(\mathrm{E};\mathrm{F})\) and \(v \in \mathcal{L}_2(\mathrm{F};\mathrm{G})\). Let \(w \in \mathcal{L}_2(\mathrm{E};\mathrm{G})\) . We obtain \(\langle uv | w \rangle = \operatorname{Tr}(v^* u^* w) = \langle v | u^* w \rangle\) whence

\[
| \langle u v | w \rangle | \leqslant \| v \| _ {2} \| u ^ {*} w \| _ {2} \leqslant \| v \| _ {2} \| u \| _ {2} \| w \|
\]

by the Cauchy--Schwarz inequality and formula (37) of EVT, V, p. 52. The result therefore follows from prop. 14.

Lemma 8. --- The map \(u \mapsto u^{*}\) from \(\mathcal{L}(\mathrm{E};\mathrm{F})\) into \(\mathcal{L}(\mathrm{F};\mathrm{E})\) defines by passage to quotients an isometric linear map from \(\mathcal{L}_{1}(\mathrm{E};\mathrm{F})\) into \(\mathcal{L}_{1}(\mathrm{F};\mathrm{E})\) .

Let \(u \in \mathcal{L}_1(\mathrm{E};\mathrm{F})\) . Since \(u\) is a Hilbert--Schmidt map, so is \(u^*\) (EVT, V, p. 54). The map \(v \mapsto v^*\) is a bijection from the set of \(v \in \mathcal{L}_2(\mathrm{E};\mathrm{F})\) such that \(\|v\| \leqslant 1\) onto the set of \(w \in \mathcal{L}_2(\mathrm{F};\mathrm{E})\) such that \(\|w\| \leqslant 1\) ; for every \(v \in \mathcal{L}_2(\mathrm{E};\mathrm{F})\)

with \(\|v\|\leqslant1\) , we have \(\langle v|u\rangle=\langle u^{*}|v^{*}\rangle\) (EVT, V, p. 54, formula (42)), whence the result by prop. 14.

PROPOSITION 16. --- Let \(E_{1}\) and \(F_{1}\) be Hilbert spaces. Let \(u\) be in \(\mathcal{L}_{1}(\mathrm{E};\mathrm{F})\), \(v\) be in \(\mathcal{L}(\mathrm{E}_{1};\mathrm{E})\) and \(v_{2}\) be in \(\mathcal{L}(\mathrm{F};\mathrm{F}_{1})\). We then have \(v_{2}uv_{1} \in \mathcal{L}_{1}(\mathrm{E}_{1};\mathrm{F}_{1})\) and \(\|v_{2}uv_{1}\|_{1} \leqslant \|v_{2}\|\|v_{1}\|\|u\|_{1}\) .

Let \(w \in \mathcal{L}_2(\mathrm{E};\mathrm{F}_1)\) such that \(\| w\| \leqslant 1\) . We have \(v_{2}u \in \mathcal{L}_{2}(\mathrm{E};\mathrm{F}_{1})\) (EVT, V, p. 52, formula (36)). Since \(v_{2}^{*}w \in \mathcal{L}_{2}(\mathrm{E};\mathrm{F})\) (loc. cit.), it follows

\[
| \langle w | v _ {2} u \rangle | = | \langle v _ {2} ^ {*} w | u \rangle | \leqslant \| v _ {2} \| \| u \| _ {1}
\]

(formula (11)), whence \(v_2u \in \mathcal{L}_1(\mathrm{E};\mathrm{F}_1)\) and \(\| v_2u\| _1 \leqslant \| v_2\| \| u\| _1\) (prop. 14). Let \(v_{1} \in \mathcal{L}(\mathrm{E}_{1};\mathrm{E})\) ; since \(uv_{1} = (v_{1}^{*}u^{*})^{*}\) , we have \(uv_{1} \in \mathcal{L}_{1}(\mathrm{E}_{1};\mathrm{F})\) and \(\| uv_{1}\|_{1} \leqslant \| v_{1}^{*}\| \| u^{*}\|_{1} = \| v_{1}\| \| u\|_{1}\) (lemma 8).

The proposition follows immediately from these inequalities.

PROPOSITION 17. — The space \(\mathcal{L}_{1}(\mathrm{E};\mathrm{E})\) coincides with the space \(\mathcal{L}_{1}(\mathrm{E})\) of endomorphisms of finite trace of E, and we have \(|\operatorname{Tr}(u)|\leqslant \| u\|_{1}\) for every endomorphism \(u\) of E with finite trace.

We may assume that K = C. The space \(\mathcal{L}_{1}(E;E)\) contains by definition the set of positive endomorphisms of finite trace, and hence \(\mathcal{L}_{1}(E)\subset\mathcal{L}_{1}(E;E)\) (EVT, p. 50, def. 8). Conversely, let us show that \(\mathcal{L}_{1}(E;E)\) is contained in \(\mathcal{L}_{1}(E)\) .

Let \(u \in \mathcal{L}_{1}(\mathrm{E};\mathrm{E})\) . We have \(u^{*} \in \mathcal{L}_{1}(\mathrm{E};\mathrm{E})\) (lemma 8); it therefore suffices to show that the Hermitian elements of \(\mathcal{L}_{1}(\mathrm{E};\mathrm{E})\) are of finite trace (lemma 2 of I, p. 96).

Let u be such an endomorphism. It is compact (remark 2, p. 169). Let B be an orthonormal basis of E such that u is diagonal in the basis B (th. 1 of IV, p. 149), and let \(\lambda\) be the family of eigenvalues of u relative to B. The endomorphism \(|u|\) is diagonalizable in the basis B and the family of its eigenvalues is \(|\lambda|\) (prop. 1, d) of IV, p. 147). Since \(|u|\) is of finite trace, this latter family is summable (prop. 12 of IV, p. 164), hence the family \(\lambda\) is summable. Consequently, u is of finite trace (loc. cit.), and we have \(|\mathrm{Tr}(u)| \leqslant \mathrm{Tr}(|u|) = \|u\|_{1}\) .

Consequently, the space \(\mathcal{L}_{1}(\mathrm{E})\) is a self-adjoint two-sided ideal of the C*-algebra \(\mathcal{L}(\mathrm{E})\) . If E is of infinite dimension, this ideal is not closed in \(\mathcal{L}(\mathrm{E})\) .

Propositions 17 and 16 show that the map \(b\) from \(\mathcal{L}_1(\mathrm{E};\mathrm{F})\times \mathcal{L}(\mathrm{F};\mathrm{E})\) into \(\mathbf{C}\) defined by \(b(u,v) = \operatorname {Tr}(vu)\) is a continuous bilinear form which puts the spaces \(\mathcal{L}_1(\mathrm{E};\mathrm{F})\) and \(\mathcal{L}(\mathrm{F};\mathrm{E})\) in duality.

Lemma 9. --- Let \(u \in \mathcal{L}_1(\mathrm{E};\mathrm{F})\) . We have

\[
\| u\|_{1} = \sup_{\substack{v\in \mathcal{L}^{c}(\mathrm{F};\mathrm{E})\\ \| v\| \leqslant 1}}|b(u,v)|.
\]

Since every Hilbert--Schmidt map is compact (cor. 2 of IV, p. 165), we have

\[
\| u\|_{1}\leqslant \sup_{\substack{v\in \mathcal{L}^{\mathrm{c}}(\mathrm{E};\mathrm{F})\\ \| v\| \leqslant 1}}|\operatorname{Tr}(vu)| = \sup_{\substack{v\in \mathcal{L}^{\mathrm{c}}(\mathrm{F};\mathrm{E})\\ \| v\| \leqslant 1}}|b(u,v)|
\]

by prop. 14.

Let \(v \in \mathcal{L}^{c}(\mathrm{E};\mathrm{F})\) be of norm \(\leqslant 1\) . Since \(\mathcal{L}_{2}(\mathrm{E};\mathrm{F})\) is dense in \(\mathcal{L}^{c}(\mathrm{E};\mathrm{F})\) , there exists a sequence \((v_{n})_{n\in\mathbf{N}}\) in \(\mathcal{L}_{2}(\mathrm{E};\mathrm{F})\) which converges to v in \(\mathcal{L}(\mathrm{E};\mathrm{F})\) . The sequence \((b(u,v_{n}))_{n\in\mathbf{N}}\) converges to \(b(u,v)\) ; since \(|b(u,v_{n})| = |\operatorname{Tr}(v_{n}u)| = |\langle v_{n}^{*}|u\rangle| \leqslant \|u\|_{1}\) (prop. 14) for all \(n \in N\) , it follows that \(|b(u,v)| \leqslant \|u\|_{1}\) .

Denote by \(\theta\) the continuous linear map

\[
\theta \colon \mathcal {L} _ {1} (\mathrm{E}; \mathrm{F}) \to \mathcal {L} ^ {\mathrm{c}} (\mathrm{F}; \mathrm{E}) ^ {\prime}
\]

such that \(\theta(u)(v) = b(u,v)\) .

PROPOSITION 18. --- The map θ is an isometric isomorphism.

By lemma 9, the map \(\theta\) is isometric, and it suffices to show that \(\theta\) is surjective.

Let \(\lambda \in \mathcal{L}^{\mathrm{c}}(\mathrm{F};\mathrm{E})'\) . Since \(\| v\| \leqslant \| v\|_{2}\) for all \(v\in \mathcal{L}_2(\mathrm{F};\mathrm{E})\) (EVT, V, p. 52, formula (33)), the restriction of \(\lambda\) to the subspace \(\mathcal{L}_2(\mathrm{F};\mathrm{E})\) is a continuous linear form on the Hilbert space \(\mathcal{L}_2(\mathrm{F};\mathrm{E})\) . There therefore exists an element \(u\) of \(\mathcal{L}_2(\mathrm{F};\mathrm{E})\) such that \(\lambda (v) = \langle u|v\rangle\) for all \(v\in \mathcal{L}_2(\mathrm{F};\mathrm{E})\) (EVT, V, p. 15, th. 3).

For every \(w \in \mathcal{L}_2(\mathrm{E};\mathrm{F})\) of norm \(\leqslant 1\) , we have \(|\langle w|u\rangle| = |\lambda(w)| \leqslant \| \lambda \|\) . Consequently, \(u\) is a nuclear map from \(\mathrm{E}\) into \(\mathrm{F}\) (prop. 14).

The continuous linear forms \(\lambda\) and \(\theta(u)\) are equal on the dense subspace \(\mathcal{L}_{2}(\mathrm{F},\mathrm{E})\) of \(\mathcal{L}^{\mathrm{c}}(\mathrm{F};\mathrm{E})\) . It follows that \(\lambda=\theta(u)\) , which concludes the proof.

COROLLARY. --- The normed space \(\mathcal{L}_{1}(\mathrm{E};\mathrm{F})\) is a Banach space.

The assertion follows from the proposition and from EVT, III, p. 24, cor. 2, since the space \(\mathcal{L}^{\mathrm{c}}(\mathrm{F};\mathrm{E})\) is a normed space.

